% SEGMENT OLP-0542-S01
% Part: set-theory
% Chapter: z
% Section: powerset

\documentclass[../../../include/open-logic-section]{subfiles}

\begin{document}

\olfileid{sth}{z}{power}
\olsection{அடுக்குக்கணங்கள்}

மேலும் ஓர் அடிகோளுடன் தொடர்வோம்:

\begin{axiom}[அடுக்குக்கணங்கள்]
எந்தக் கணம் $A$ க்கும், $\Pow{A} = \Setabs{x}{x \subseteq A}$ என்ற
கணம் உள்ளது.
\[
	\forall A \exists P \forall x(x \in P \liff (\forall z \in x)z \in A)
\]
\end{axiom}

இதற்கான நமது நியாயம் மிகவும் நேரடியானது. படிநிலை $S$ இல் $A$
உருவாகிறது என்க. அப்போது $A$ இன் அனைத்து உறுப்புகளும் $S$ க்கு முன்பு
கிடைத்தன (\stagesacc{} இன்படி). ஆகவே பிரிப்பிற்கான நமது
நியாயப்படுத்தலில் செய்ததைப்போல் பகுத்தறிந்தால், $A$ இன் ஒவ்வோர் உட்கணமும்
படிநிலை $S$ ஐ அடைவதற்குள் உருவாகிறது. எனவே அவை அனைத்தும் $S$ க்குப்
பிந்திய எந்தப் படிநிலையிலும் ஒரே கணமாக உருவாக்கப்படுவதற்குக் கிடைக்கின்றன.
$S$ கடைசிப் படிநிலை அல்ல என்பதால் (\stagessucc{} இன்படி), அத்தகைய ஒரு
படிநிலை உள்ளது என்பதை அறிவோம். ஆகவே $\Pow{A}$ உள்ளது.

அடுக்குக்கணங்களின் ஓர் அழகிய விளைவு இதோ:

\begin{prop}\label{thm:Products}
எந்தக் கணங்கள் $A, B$ தரப்பட்டாலும், அவற்றின் கார்டீசியன் பெருக்கம்
$A \times B$ உள்ளது.
\end{prop}

\begin{proof}
அடுக்குக்கணங்கள் அடிகோளின்படியும்
\olref[sth][z][pairs]{prop:pairsconsequences} இன்படியும்
$\Pow{\Pow{A \cup B}}$ என்ற கணம் உள்ளது. ஆகவே பிரிப்பின்படி, பின்வரும்
கணம் உள்ளது:
\[
	C = \Setabs{z \in \Pow{\Pow{A \cup B}}}{(\exists x \in A)(\exists y \in B) z = \tuple{x, y}}.
\]
இப்போது எந்த $x \in A$, $y \in B$ க்கும்,
\olref[sth][z][pairs]{prop:pairsconsequences} இன்படி $\tuple{x, y}$ என்ற
கணம் உள்ளது. மேலும் $x, y \in A \cup B$ என்பதால்,
$\{x\}, \{x, y\} \in \Pow{A \cup B}$; ஆகவே
$\tuple{x,y} \in \Pow{\Pow{A \cup B}}$. எனவே $A \times B = C$.
\end{proof}

இந்த நிறுவலில் அடுக்குக்கணங்கள் அடிகோள் பிரிப்புடன் இடைவினை புரிகிறது.
அது வியப்புக்குரியதல்ல. பிரிப்பு இல்லாவிட்டால், அடுக்குக்கணங்கள் அடிகோள்
மிகவும் \emph{ஆற்றல்மிக்க} கொள்கையாக இருக்காது. ஏனெனில் ஒரு கணத்தின்
எந்த உட்கணங்கள் உள்ளன என்பதைப் பிரிப்பே நமக்குக் கூறுகிறது; அதன்மூலம்
ஒவ்வோர் அடுக்குக்கணமும் எவ்வளவு ``செறிவானது'' என்பதையும் அது
நிர்ணயிக்கிறது.

\begin{prob}
எந்தக் கணங்கள் $A, B$ க்கும்: (i) பொருட்களம் $A$ ஆகவும் வீச்சு $B$
ஆகவும் கொண்ட அனைத்துத் தொடர்புகளின் கணம் உள்ளது; (ii) $A$ இலிருந்து
$B$ க்குச் செல்லும் அனைத்துச் சார்புகளின் கணம் உள்ளது எனக் காட்டுக.
\end{prob}

\begin{prob}
$A$ ஒரு கணமாகவும், $\sim$ என்பது $A$ மீதான ஒரு சமானத் தொடர்பாகவும்
இருக்கட்டும். $A$ மீது $\sim$ இன் கீழ் கிடைக்கும் சமான வகுப்புகளின்
கணம், அதாவது $\equivclass{A}{\sim}$, உள்ளது என நிறுவுக.
\end{prob}

\end{document}
